%%%%%%%%%%%%%%%%%%%%%%%%%%%%%%%%%
%
%       EXERCÍCIO
%
%%%%%%%%%%%%%%%%%%%%%%%%%%%%%%%%%

\definevalorquestao

\begin{exercicioBanco}[\valorquestao]
A figura a seguir representa um quadrado com 20 cm de lado.

\begin{center}
    \begin{tikzpicture}[scale=0.4]

        % Pontos
        \coordinate (A) at (0,0);
        \coordinate (B) at (8,0);
        \coordinate (C) at (8,8);
        \coordinate (D) at (0,8);

        \coordinate (E) at (3,8);   % vértice superior
        \coordinate (F) at (0,4);   % encontro da lateral esquerda
        \coordinate (G) at (8,2);   % encontro da lateral direita

        % Região hachurada
        \fill[pattern=north east lines]
        (A)--(B)--(G)--(E)--(F)--cycle;

        % Contorno
        \draw[thick]
        (A)--(B)--(G)--(E)--(F)--cycle;

        % Contorno
        \draw[thick]
        (A)--(B)--(C)--(D)--cycle;

        % Letras
        \node[left]  at (A) {$A$};
        \node[right] at (B) {$B$};
        \node[above] at (C) {$C$};
        \node[above left] at (D) {$D$};

        % Cota x
        \draw[dashed] (D) -- ++(3,0);
        \draw[<->] (0,8.7) -- (3,8.7);
        \node at (1.5,9.15) {$x$};

        % Cota 2x
        \draw[dashed] (-0.6,8) -- (-0.6,4);
        \draw[<->] (-0.6,4) -- (-0.6,8);
        \node[left] at (-0.8,6) {$2x$};

        % Cota 6
        \draw[dashed] (8.6,8) -- (8.6,6);
        \draw[<->] (8.6,8) -- (8.6,2);
        \node[right] at (8.8,5) {$6$};

    \end{tikzpicture}
\end{center}

Determine a expressão da área y da parte hachurada da figura.

% Define as alternativas
\newcommand{\alternativas}{%
    \begin{center}
        \begin{tabularx}{\textwidth}{XXX}
            (a) \(y = x^{2} + 3x + 240\). &
            (b) \(y = -x^{2} + 3x + 280\). &
            (c) \(y = -x^{2} + 3x + 340\). \\[5pt]
            (d) \(y = x^{2} + 3x + 340\). &
            (e) \(y = -x^{2} + x + 340\).
        \end{tabularx}
    \end{center}
}

% Define a resposta correta
\newcommand{\resposta}{C}

% Lógica condicional para exibição
\ifdefstring{\atividade}{lista}{%
    \alternativas
    \vspace{0.5em}
    
    \noindent\textbf{Resposta:} letra \textbf{\resposta}.
}{%
    \ifdefstring{\modo}{objetivo}{%
        \alternativas
    }{%
        % modo = discursiva → não mostra alternativas
    }
}
\end{exercicioBanco}

